%!TEX TS-program = xelatex
%!TEX encoding = UTF-8 Unicode
\documentclass[10pt]{article}
\title{Java vs Go --- A cheatsheet}
\author{Mathieu N\'eron}
\usepackage{cheatsheet-style}

\begin{document}
\cheatsheetTitle{Java vs Go}{Mathieu N\'eron}{2026}

\noindent\small This cheatsheet anchors Go on constructs you already know in Java.
Left column: Java. Right column: Go (1.22+). Go is intentionally minimal --- the
biggest mental shifts for a Java dev are: \textbf{no inheritance}, \textbf{errors as
return values}, \textbf{structural interfaces}, \textbf{goroutines + channels} for
concurrency.

%===============================================================================
\cheatsection{Hello world \& program entry}
%===============================================================================

\begin{construct}{Program entry point}
\noindent
\begin{minipage}[t]{0.485\linewidth}
\begin{lstlisting}[style=java, title=Java]
public class Hello {
  public static void main(String[] args) {
    System.out.println("Hello, " + args[0]);
  }
}
// javac Hello.java && java Hello world
\end{lstlisting}
\end{minipage}\hfill
\begin{minipage}[t]{0.485\linewidth}
\begin{lstlisting}[style=go, title=Go]
package main

import (
    "fmt"
    "os"
)

func main() {
    fmt.Println("Hello,", os.Args[1])
}
// go run hello.go world  /  go build && ./hello world
\end{lstlisting}
\end{minipage}
\end{construct}

\begin{construct}{File / module / package layout}
\noindent
\begin{minipage}[t]{0.485\linewidth}
\begin{lstlisting}[style=java, title=Java]
// One public class per file.
// src/main/java/com/acme/Hello.java
package com.acme;
// pom.xml or build.gradle for deps
\end{lstlisting}
\end{minipage}\hfill
\begin{minipage}[t]{0.485\linewidth}
\begin{lstlisting}[style=go, title=Go]
// Each directory == one package. All .go files in
// the dir share the package namespace.
// go.mod at module root:
module github.com/acme/svc
go 1.22
require github.com/foo/bar v1.2.3
\end{lstlisting}
\end{minipage}
\end{construct}

%===============================================================================
\cheatsection{Variables, constants, type inference}
%===============================================================================

\begin{construct}{Local variables}
\noindent
\begin{minipage}[t]{0.485\linewidth}
\begin{lstlisting}[style=java, title=Java]
int x = 5;
final String s = "hi";
var list = new ArrayList<Integer>();   // Java 10+
\end{lstlisting}
\end{minipage}\hfill
\begin{minipage}[t]{0.485\linewidth}
\begin{lstlisting}[style=go, title=Go]
var x int = 5
var s = "hi"          // type inferred
y := 5                // short decl, only inside funcs
const Pi = 3.14159    // typed by use site
const (
    A = iota          // 0
    B                 // 1
    C                 // 2
)
\end{lstlisting}
\end{minipage}
\end{construct}

\begin{construct}{Zero values \& visibility}
\noindent
\begin{minipage}[t]{0.485\linewidth}
\begin{lstlisting}[style=java, title=Java]
// Fields default: 0, false, null. Locals must be init'd.
// Visibility via keywords: public/protected/private.
public class C { private int x; }
\end{lstlisting}
\end{minipage}\hfill
\begin{minipage}[t]{0.485\linewidth}
\begin{lstlisting}[style=go, title=Go]
// All declared vars get the type's zero value
// (0, "", false, nil). No null-pointer surprises on init.
// Visibility: Capitalized = exported, lowercase = unexported.
type C struct { x int }   // x unexported (private)
type D struct { X int }   // X exported   (public)
\end{lstlisting}
\end{minipage}
\end{construct}

%===============================================================================
\cheatsection{Primitive types \& conversions}
%===============================================================================

\begin{construct}{Numeric types}
\noindent
\begin{minipage}[t]{0.485\linewidth}
\begin{lstlisting}[style=java, title=Java]
byte b; short s; int i; long l;
float f; double d;
int n = Integer.parseInt("42");
double x = Double.parseDouble("1.5");
// Implicit widening: int -> long. Narrowing requires cast.
long L = i;    int j = (int) L;
\end{lstlisting}
\end{minipage}\hfill
\begin{minipage}[t]{0.485\linewidth}
\begin{lstlisting}[style=go, title=Go]
// int8/16/32/64, uint*, float32/64, complex64/128
// `int`/`uint` are platform-width (usually 64).
import "strconv"
n, err := strconv.Atoi("42")
x, err := strconv.ParseFloat("1.5", 64)
// NO implicit conversions, even widening:
var i int32 = 5
var L int64 = int64(i)        // explicit
\end{lstlisting}
\end{minipage}
\end{construct}

\begin{construct}{Bool, byte, rune, string}
\noindent
\begin{minipage}[t]{0.485\linewidth}
\begin{lstlisting}[style=java, title=Java]
boolean b = true;
char c = 'A';                // 16-bit UTF-16 unit
byte[] bs = "hi".getBytes(StandardCharsets.UTF_8);
String s = new String(bs, StandardCharsets.UTF_8);
\end{lstlisting}
\end{minipage}\hfill
\begin{minipage}[t]{0.485\linewidth}
\begin{lstlisting}[style=go, title=Go]
b := true
c := 'A'                     // rune == int32 (Unicode pt)
bs := []byte("hi")           // UTF-8 bytes
s := string(bs)
// Strings are immutable byte sequences (typically UTF-8).
// for i, r := range s iterates runes, not bytes.
\end{lstlisting}
\end{minipage}
\end{construct}

%===============================================================================
\cheatsection{Strings}
%===============================================================================

\begin{construct}{Interpolation \& formatting}
\noindent
\begin{minipage}[t]{0.485\linewidth}
\begin{lstlisting}[style=java, title=Java]
String name = "world";
String s1 = "Hello, " + name;
String s2 = String.format("%s = %.2f", "pi", 3.14);
// Java 15+ text block:
String body = """
    Hello,
    World
    """;
\end{lstlisting}
\end{minipage}\hfill
\begin{minipage}[t]{0.485\linewidth}
\begin{lstlisting}[style=go, title=Go]
name := "world"
s1 := "Hello, " + name
s2 := fmt.Sprintf("%s = %.2f", "pi", 3.14)
// Raw string literal (backticks): keeps newlines, no escapes
body := `
    Hello,
    World
    `
\end{lstlisting}
\end{minipage}
\end{construct}

\begin{construct}{Common operations}
\noindent
\begin{minipage}[t]{0.485\linewidth}
\begin{lstlisting}[style=java, title=Java]
s.length(); s.toUpperCase(); s.trim();
s.startsWith("H"); s.contains("ll");
String[] parts = s.split(",");
String joined = String.join("-", parts);
\end{lstlisting}
\end{minipage}\hfill
\begin{minipage}[t]{0.485\linewidth}
\begin{lstlisting}[style=go, title=Go]
import "strings"
len(s)                                // BYTES (not runes)
strings.ToUpper(s); strings.TrimSpace(s)
strings.HasPrefix(s, "H"); strings.Contains(s, "ll")
parts := strings.Split(s, ",")
joined := strings.Join(parts, "-")
\end{lstlisting}
\end{minipage}
\end{construct}

%===============================================================================
\cheatsection{Control flow}
%===============================================================================

\begin{construct}{If / else (no ternary)}
\noindent
\begin{minipage}[t]{0.485\linewidth}
\begin{lstlisting}[style=java, title=Java]
if (x > 0) { ... }
else if (x == 0) { ... }
else { ... }

int sign = x > 0 ? 1 : (x < 0 ? -1 : 0);
\end{lstlisting}
\end{minipage}\hfill
\begin{minipage}[t]{0.485\linewidth}
\begin{lstlisting}[style=go, title=Go]
if x > 0 {
    ...
} else if x == 0 {
    ...
} else { ... }

// No ternary. Idiom: a tiny helper or explicit if.
sign := 0
if x > 0 { sign = 1 } else if x < 0 { sign = -1 }

// "if with init" is a Go idiom:
if v, err := f(); err == nil { use(v) }
\end{lstlisting}
\end{minipage}
\end{construct}

\begin{construct}{Switch}
\noindent
\begin{minipage}[t]{0.485\linewidth}
\begin{lstlisting}[style=java, title=Java]
String label = switch (x) {
  case 1, 2  -> "small";
  case 3     -> "three";
  default    -> "other";
};
// Pattern switch (21+):
switch (s) {
  case Circle c -> ...;
  case null     -> ...;
}
\end{lstlisting}
\end{minipage}\hfill
\begin{minipage}[t]{0.485\linewidth}
\begin{lstlisting}[style=go, title=Go]
switch x {
case 1, 2:
    label = "small"
case 3:
    label = "three"
default:
    label = "other"
}
// No implicit fallthrough; use `fallthrough` keyword.
// "Type switch":
switch v := s.(type) {
case Circle: ...
case nil:    ...
}
\end{lstlisting}
\end{minipage}
\end{construct}

%===============================================================================
\cheatsection{Loops (just `for`)}
%===============================================================================

\begin{construct}{All loops use `for`}
\noindent
\begin{minipage}[t]{0.485\linewidth}
\begin{lstlisting}[style=java, title=Java]
for (int i = 0; i < 10; i++) { ... }
while (cond) { ... }
do { ... } while (cond);
for (var x : xs) { ... }
\end{lstlisting}
\end{minipage}\hfill
\begin{minipage}[t]{0.485\linewidth}
\begin{lstlisting}[style=go, title=Go]
for i := 0; i < 10; i++ { ... }   // classic
for cond { ... }                  // == while
for { ... }                       // == while(true)
for i, x := range xs { ... }      // index, value
for k, v := range m { ... }       // map
for _, x := range xs { ... }      // discard index
// 1.22+: each loop iter creates fresh `i` (no aliasing trap)
\end{lstlisting}
\end{minipage}
\end{construct}

%===============================================================================
\cheatsection{Collections \& data structures}
%===============================================================================

\begin{construct}{Arrays vs slices vs maps}
\noindent
\begin{minipage}[t]{0.485\linewidth}
\begin{lstlisting}[style=java, title=Java]
int[] a = {1, 2, 3};                 // fixed
List<Integer> xs = new ArrayList<>(List.of(1,2,3));
xs.add(4); xs.get(0); xs.size();
Set<String> set = new HashSet<>(List.of("a","b"));
Map<String,Integer> m = new HashMap<>();
m.put("a", 1);
\end{lstlisting}
\end{minipage}\hfill
\begin{minipage}[t]{0.485\linewidth}
\begin{lstlisting}[style=go, title=Go]
a := [3]int{1, 2, 3}                 // ARRAY (fixed size)
xs := []int{1, 2, 3}                 // SLICE (dynamic view)
xs = append(xs, 4); xs[0]; len(xs); cap(xs)
xs[1:3]                              // sub-slice (no copy!)

// No built-in set; use map[T]struct{}
set := map[string]struct{}{"a": {}, "b": {}}
_, ok := set["a"]                    // membership

m := map[string]int{"a": 1}
m["a"] = 1; v, ok := m["a"]; delete(m, "a")
\end{lstlisting}
\end{minipage}
\end{construct}

\begin{construct}{Slice gotcha}
\noindent
\begin{minipage}[t]{0.485\linewidth}
\begin{lstlisting}[style=java, title=Java]
// subList returns a view; modifying it mutates the parent.
// Same general "view vs copy" hazard, less surprising
// because typical code uses fresh ArrayList copies.
\end{lstlisting}
\end{minipage}\hfill
\begin{minipage}[t]{0.485\linewidth}
\begin{lstlisting}[style=go, title=Go]
// Slices share backing array. Sub-slicing is a VIEW.
a := []int{1,2,3,4,5}
b := a[1:3]                  // {2,3}, shares storage
b[0] = 99                    // mutates a[1] too!
// To fully detach: copy
c := make([]int, len(b)); copy(c, b)
\end{lstlisting}
\end{minipage}
\end{construct}

%===============================================================================
\cheatsection{Functions}
%===============================================================================

\begin{construct}{Definition, varargs, multiple returns}
\noindent
\begin{minipage}[t]{0.485\linewidth}
\begin{lstlisting}[style=java, title=Java]
int add(int a, int b) { return a + b; }
int sum(int... xs) {
    int t = 0; for (var x : xs) t += x; return t;
}
sum(1, 2, 3);
// No multi-return: use record or out-params.
\end{lstlisting}
\end{minipage}\hfill
\begin{minipage}[t]{0.485\linewidth}
\begin{lstlisting}[style=go, title=Go]
func add(a, b int) int { return a + b }

func sum(xs ...int) int {
    t := 0
    for _, x := range xs { t += x }
    return t
}

// Multi-return is idiomatic, especially (T, error):
func divmod(x, y int) (int, int) { return x/y, x%y }
q, r := divmod(10, 3)
\end{lstlisting}
\end{minipage}
\end{construct}

\begin{construct}{Lambdas / closures}
\noindent
\begin{minipage}[t]{0.485\linewidth}
\begin{lstlisting}[style=java, title=Java]
Function<Integer,Integer> sq = x -> x * x;
list.forEach(System.out::println);
\end{lstlisting}
\end{minipage}\hfill
\begin{minipage}[t]{0.485\linewidth}
\begin{lstlisting}[style=go, title=Go]
sq := func(x int) int { return x * x }
adder := func() func(int) int {
    sum := 0
    return func(x int) int { sum += x; return sum }
}()
adder(2); adder(3)            // 2, 5 (closure captures sum)
\end{lstlisting}
\end{minipage}
\end{construct}

%===============================================================================
\cheatsection{Structs \& interfaces (NO inheritance)}
%===============================================================================

\begin{construct}{Class vs struct + methods}
\noindent
\begin{minipage}[t]{0.485\linewidth}
\begin{lstlisting}[style=java, title=Java]
public class Point {
  private final double x, y;
  public Point(double x, double y) {
    this.x = x; this.y = y;
  }
  public double dist(Point o) {
    return Math.hypot(x - o.x, y - o.y);
  }
}
\end{lstlisting}
\end{minipage}\hfill
\begin{minipage}[t]{0.485\linewidth}
\begin{lstlisting}[style=go, title=Go]
type Point struct {
    X, Y float64               // exported fields
}

// Method = func with a receiver
func (p Point) Dist(o Point) float64 {
    dx, dy := p.X-o.X, p.Y-o.Y
    return math.Sqrt(dx*dx + dy*dy)
}

// Constructor: just a function, no `new` keyword needed
func NewPoint(x, y float64) Point { return Point{X: x, Y: y} }
\end{lstlisting}
\end{minipage}
\end{construct}

\begin{construct}{Pointer vs value receiver}
\noindent
\begin{minipage}[t]{0.485\linewidth}
\begin{lstlisting}[style=java, title=Java]
// Java: objects always passed by reference (the reference
// is by value). No equivalent distinction.
\end{lstlisting}
\end{minipage}\hfill
\begin{minipage}[t]{0.485\linewidth}
\begin{lstlisting}[style=go, title=Go]
type Counter struct { n int }
func (c Counter) GetVal() int { return c.n }   // value rcv
func (c *Counter) Inc()       { c.n++ }        // ptr rcv
// Use pointer receivers when method mutates state OR
// the struct is large. Be consistent within a type.
\end{lstlisting}
\end{minipage}
\end{construct}

\begin{construct}{Interfaces (structural, implicit)}
\noindent
\begin{minipage}[t]{0.485\linewidth}
\begin{lstlisting}[style=java, title=Java]
public interface Stringer { String toStr(); }
public class Foo implements Stringer {
  public String toStr() { return "foo"; }
}
// Nominal: must declare `implements`.
\end{lstlisting}
\end{minipage}\hfill
\begin{minipage}[t]{0.485\linewidth}
\begin{lstlisting}[style=go, title=Go]
type Stringer interface {
    String() string
}

type Foo struct{}
func (f Foo) String() string { return "foo" }
// Foo SATISFIES Stringer automatically. No "implements".
// Interfaces are STRUCTURAL: shape, not name, matches.
\end{lstlisting}
\end{minipage}
\end{construct}

\begin{construct}{Composition replaces inheritance}
\noindent
\begin{minipage}[t]{0.485\linewidth}
\begin{lstlisting}[style=java, title=Java]
class Animal { void breathe() { ... } }
class Dog extends Animal { void bark() { ... } }
Dog d = new Dog(); d.breathe(); d.bark();
\end{lstlisting}
\end{minipage}\hfill
\begin{minipage}[t]{0.485\linewidth}
\begin{lstlisting}[style=go, title=Go]
type Animal struct{}
func (a Animal) Breathe() { ... }

type Dog struct {
    Animal                // embedded type (anonymous field)
}
func (d Dog) Bark() { ... }

d := Dog{}
d.Breathe()              // promoted from embedded Animal
d.Bark()
// No virtual dispatch. Method set "promoted" via embedding.
\end{lstlisting}
\end{minipage}
\end{construct}

%===============================================================================
\cheatsection{Generics (Go 1.18+)}
%===============================================================================

\begin{construct}{Generic function \& type}
\noindent
\begin{minipage}[t]{0.485\linewidth}
\begin{lstlisting}[style=java, title=Java]
<T extends Comparable<T>> T max(T a, T b) {
  return a.compareTo(b) >= 0 ? a : b;
}
class Box<T> { final T v; Box(T v) { this.v = v; } }
\end{lstlisting}
\end{minipage}\hfill
\begin{minipage}[t]{0.485\linewidth}
\begin{lstlisting}[style=go, title=Go]
func Max[T cmp.Ordered](a, b T) T {
    if a >= b { return a }
    return b
}

type Box[T any] struct { V T }
b := Box[int]{V: 5}
\end{lstlisting}
\end{minipage}
\end{construct}

\begin{construct}{Constraints}
\noindent
\begin{minipage}[t]{0.485\linewidth}
\begin{lstlisting}[style=java, title=Java]
// Constraint = upper bound on a class hierarchy.
<T extends Number & Comparable<T>>
T sum(List<T> xs) { ... }
\end{lstlisting}
\end{minipage}\hfill
\begin{minipage}[t]{0.485\linewidth}
\begin{lstlisting}[style=go, title=Go]
// Constraints are interfaces (possibly with type sets):
type Number interface {
    ~int | ~int64 | ~float32 | ~float64
}
func Sum[T Number](xs []T) T {
    var t T
    for _, x := range xs { t += x }
    return t
}
\end{lstlisting}
\end{minipage}
\end{construct}

%===============================================================================
\cheatsection{Error handling (no exceptions)}
%===============================================================================

\begin{construct}{Errors are values}
\noindent
\begin{minipage}[t]{0.485\linewidth}
\begin{lstlisting}[style=java, title=Java]
try (var f = Files.newBufferedReader(p)) {
    return f.readLine();
} catch (IOException e) {
    log.error("read failed", e);
    throw new MyException("oops", e);
}
\end{lstlisting}
\end{minipage}\hfill
\begin{minipage}[t]{0.485\linewidth}
\begin{lstlisting}[style=go, title=Go]
data, err := os.ReadFile(p)
if err != nil {
    return fmt.Errorf("read failed: %w", err)  // wrap
}
// No try/catch. Functions return (T, error).
// Caller MUST check err. Idiom: "if err != nil { return }"
\end{lstlisting}
\end{minipage}
\end{construct}

\begin{construct}{Custom error \& sentinel checks}
\noindent
\begin{minipage}[t]{0.485\linewidth}
\begin{lstlisting}[style=java, title=Java]
public class NotFound extends RuntimeException { ... }
catch (NotFound e) { ... }
catch (RuntimeException e) { ... }
\end{lstlisting}
\end{minipage}\hfill
\begin{minipage}[t]{0.485\linewidth}
\begin{lstlisting}[style=go, title=Go]
var ErrNotFound = errors.New("not found")
type ValidationError struct{ Field string }
func (e *ValidationError) Error() string { ... }

if errors.Is(err, ErrNotFound) { ... }
var ve *ValidationError
if errors.As(err, &ve) { use(ve.Field) }
\end{lstlisting}
\end{minipage}
\end{construct}

\begin{construct}{Cleanup: defer (replaces try-with-resources)}
\noindent
\begin{minipage}[t]{0.485\linewidth}
\begin{lstlisting}[style=java, title=Java]
try (var f = new FileInputStream(p)) {
    // f auto-closed on exit
}
\end{lstlisting}
\end{minipage}\hfill
\begin{minipage}[t]{0.485\linewidth}
\begin{lstlisting}[style=go, title=Go]
f, err := os.Open(p)
if err != nil { return err }
defer f.Close()           // runs when func returns
// Multiple defers run in LIFO order. Args evaluated NOW.
\end{lstlisting}
\end{minipage}
\end{construct}

\begin{construct}{Panic / recover (rare)}
\noindent
\begin{minipage}[t]{0.485\linewidth}
\begin{lstlisting}[style=java, title=Java]
// throw new RuntimeException("..."); is the analog,
// but Java idiom uses exceptions widely.
\end{lstlisting}
\end{minipage}\hfill
\begin{minipage}[t]{0.485\linewidth}
\begin{lstlisting}[style=go, title=Go]
panic("unrecoverable")               // nuke goroutine
defer func() {
    if r := recover(); r != nil { ... }
}()
// Reserve panic for truly unrecoverable bugs.
// Normal errors -> return error.
\end{lstlisting}
\end{minipage}
\end{construct}

%===============================================================================
\cheatsection{Modules / packages / imports}
%===============================================================================

\begin{construct}{Imports}
\noindent
\begin{minipage}[t]{0.485\linewidth}
\begin{lstlisting}[style=java, title=Java]
package com.acme.svc;
import java.util.List;
import java.util.stream.Collectors;
// Maven coords in pom.xml. Repackaged into uberjar.
\end{lstlisting}
\end{minipage}\hfill
\begin{minipage}[t]{0.485\linewidth}
\begin{lstlisting}[style=go, title=Go]
package svc

import (
    "fmt"
    "time"
    "github.com/acme/lib"           // module path
    log "github.com/sirupsen/logrus" // alias
)
// go mod tidy resolves go.mod / go.sum
\end{lstlisting}
\end{minipage}
\end{construct}

%===============================================================================
\cheatsection{Concurrency (goroutines + channels)}
%===============================================================================

\begin{construct}{Goroutines vs threads}
\noindent
\begin{minipage}[t]{0.485\linewidth}
\begin{lstlisting}[style=java, title=Java]
Thread t = new Thread(() -> work());
t.start(); t.join();

ExecutorService pool = Executors.newFixedThreadPool(4);
Future<Integer> f = pool.submit(() -> compute());
int v = f.get();
\end{lstlisting}
\end{minipage}\hfill
\begin{minipage}[t]{0.485\linewidth}
\begin{lstlisting}[style=go, title=Go]
go work()                   // fire-and-forget goroutine
// Goroutines are cheap (~few KB stack). Spawn millions.
// No "join" — coordinate via channels or sync.WaitGroup.

var wg sync.WaitGroup
for _, item := range items {
    wg.Add(1)
    go func(x Item) { defer wg.Done(); process(x) }(item)
}
wg.Wait()
\end{lstlisting}
\end{minipage}
\end{construct}

\begin{construct}{Channels}
\noindent
\begin{minipage}[t]{0.485\linewidth}
\begin{lstlisting}[style=java, title=Java]
BlockingQueue<Integer> q = new ArrayBlockingQueue<>(8);
q.put(42);
int v = q.take();
\end{lstlisting}
\end{minipage}\hfill
\begin{minipage}[t]{0.485\linewidth}
\begin{lstlisting}[style=go, title=Go]
ch := make(chan int, 8)     // buffered (8) channel
ch <- 42                    // send
v := <-ch                   // receive
close(ch)                   // close (only sender)
for v := range ch { ... }   // drain until closed

// "Don't communicate by sharing memory; share memory by
//  communicating." — channels are the primary primitive.
\end{lstlisting}
\end{minipage}
\end{construct}

\begin{construct}{Select (multiplex)}
\noindent
\begin{minipage}[t]{0.485\linewidth}
\begin{lstlisting}[style=java, title=Java]
// Closest analog: CompletableFuture.anyOf(...).
CompletableFuture.anyOf(f1, f2)
    .thenAccept(System.out::println);
\end{lstlisting}
\end{minipage}\hfill
\begin{minipage}[t]{0.485\linewidth}
\begin{lstlisting}[style=go, title=Go]
select {
case v := <-ch1:
    use(v)
case v := <-ch2:
    use(v)
case <-time.After(2 * time.Second):
    return errors.New("timeout")
case <-ctx.Done():
    return ctx.Err()
}
\end{lstlisting}
\end{minipage}
\end{construct}

\begin{construct}{Context (cancel + deadline)}
\noindent
\begin{minipage}[t]{0.485\linewidth}
\begin{lstlisting}[style=java, title=Java]
// Roughly: Future.cancel(true), interrupt flag, or
// scheduled-future timeouts. No standardized propagation.
\end{lstlisting}
\end{minipage}\hfill
\begin{minipage}[t]{0.485\linewidth}
\begin{lstlisting}[style=go, title=Go]
ctx, cancel := context.WithTimeout(parent, 5*time.Second)
defer cancel()
result, err := callRPC(ctx, req)
// Context flows through the call tree as the first arg.
// Idiom: func F(ctx context.Context, ...) (...).
\end{lstlisting}
\end{minipage}
\end{construct}

\begin{construct}{Mutex / atomic}
\noindent
\begin{minipage}[t]{0.485\linewidth}
\begin{lstlisting}[style=java, title=Java]
synchronized (lock) { ... }
ReentrantLock lk = new ReentrantLock();
AtomicInteger ai = new AtomicInteger();
ai.incrementAndGet();
\end{lstlisting}
\end{minipage}\hfill
\begin{minipage}[t]{0.485\linewidth}
\begin{lstlisting}[style=go, title=Go]
var mu sync.Mutex
mu.Lock(); defer mu.Unlock()
// Run with `go test -race` to catch data races.

import "sync/atomic"
var n atomic.Int64
n.Add(1); v := n.Load()
\end{lstlisting}
\end{minipage}
\end{construct}

%===============================================================================
\cheatsection{Idiom appendix}
%===============================================================================

\begin{construct}{nil / equality / common gotchas}
\noindent
\begin{minipage}[t]{\linewidth}
\small
\begin{itemize}[leftmargin=*, itemsep=1pt, topsep=1pt]
  \item \textbf{nil is typed.} A \code{var p *T} is nil; an interface value is nil only if both its type AND value are nil. Returning \code{(*MyErr)(nil)} as \code{error} is NOT nil --- common bug.
  \item \textbf{Equality.} \code{==} compares structs field-by-field, but only if all fields are comparable (no slices/maps/funcs). Slices/maps/funcs are not \code{==}-comparable; use \code{reflect.DeepEqual} or hand-roll.
  \item \textbf{Range copies.} \code{for _, x := range xs} gives \code{x} as a copy. Take \code{\&xs[i]} or use index for mutation. (Pre-1.22 also: same loop variable address reused; fixed in 1.22.)
  \item \textbf{No constructors.} Just write a \code{NewT(...)} function returning \code{T} or \code{*T}.
  \item \textbf{Errors are values, not control flow.} Wrap with \code{fmt.Errorf("ctx: \%w", err)}; unwrap with \code{errors.Is/As}. Never use panic for normal flow.
  \item \textbf{Goroutine leaks.} A goroutine that blocks on a never-closed channel never exits. Always have a cancellation path (\code{ctx} or close-the-channel).
  \item \textbf{Capitalization == visibility.} \code{Foo} exported, \code{foo} package-private. JSON tags use struct tags: \code{`json:"foo,omitempty"`}.
  \item \textbf{Tooling.} \code{go fmt} (or \code{gofmt}) is non-negotiable; \code{go vet}, \code{staticcheck}, \code{go test}, \code{go test -race}, \code{go build}, \code{go mod tidy}.
\end{itemize}
\end{minipage}
\end{construct}

\end{document}
